\input{configTD}

% ── Poly à trous : commandes spécifiques ──────────────────────────

\newcommand{\atrou}[2][3cm]{%
  \ifthenelse{\boolean{showSolutions}}{%
    \par\vspace{0.5em}%
    \begin{mdframed}[backgroundcolor=ThemeLight, linewidth=0pt,
      skipabove=0pt, skipbelow=0pt,
      innertopmargin=8pt, innerbottommargin=8pt]%
    #2%
    \end{mdframed}%
    \par\vspace{0.5em}%
  }{%
    \par\vspace{0.5em}%
    \noindent\fbox{\begin{minipage}[c][#1]%
      {\dimexpr\textwidth-2\fboxsep-2\fboxrule\relax}%
      \color{white}\rule{0pt}{0pt}%
    \end{minipage}}%
    \par\vspace{0.5em}%
  }%
}

\newcommand{\val}[1]{%
  \ifthenelse{\boolean{showSolutions}}{#1}{\phantom{#1}}%
}

\title{\sffamily\bfseries Probabilités --- CM1\\[0.3em]
  \Large Espérance et variance des variables discrètes}
\author{}
\date{30 mars 2026}

\begin{document}
\sffamily
\maketitle
\thispagestyle{empty}

% ====================================================================
\section*{1. Le problème}
% ====================================================================

Commençons par un problème qui n'a rien à voir
--- a priori --- avec les probabilités.

\begin{center}
\begin{mdframed}[backgroundcolor=ThemeLight, linecolor=Theme,
  linewidth=1.5pt, innertopmargin=10pt, innerbottommargin=10pt]
\centering\large\itshape
\og Trois charges ponctuelles sont posées sur une poutre.\\[0.3em]
Où placer l'appui pour que le système soit en équilibre\,?\fg{}
\end{mdframed}
\end{center}

Considérons une poutre rigide sans masse, graduée en mètres,
sur laquelle on place trois charges :

\begin{center}
\renewcommand{\arraystretch}{1.3}
\begin{tabular}{lccc}
  \hline
  \textbf{Charge} & $A$ & $B$ & $C$ \\
  \hline
  Position $x_i$ (m) & $1$ & $3$ & $7$ \\
  Masse $m_i$ (kg)   & $2$ & $5$ & $3$ \\
  \hline
\end{tabular}
\end{center}

\textbf{Question.}
Où faut-il placer l'appui pour que la poutre soit en équilibre ?

\atrou[3.6cm]{%
On calcule le \textbf{centre de gravité},
c'est-à-dire la moyenne des positions pondérée par les masses :
\[
  \bar{x}
  = \frac{m_A \, x_A + m_B \, x_B + m_C \, x_C}
         {m_A + m_B + m_C}
  = \frac{2 \times 1 + 5 \times 3 + 3 \times 7}{2+5+3}
  = \frac{38}{10}
  = 3{,}8\;\text{m}.
\]
}

\textbf{Question.}
La formule se simplifie si la masse totale vaut~$1$.
Que devient-elle ?

\atrou[2cm]{%
\[
  \bar{x} = \sum_{i} m_i \, x_i.
\]
}

% ====================================================================
\section*{2. Du centre de gravité à l'espérance}
% ====================================================================

Imaginons maintenant que les \og masses \fg{} ne soient plus
des kilogrammes, mais des \textbf{probabilités} :
elles sont positives et leur somme vaut~$1$.

\medskip

Soit $X$ une variable aléatoire discrète prenant les valeurs
$x_1, x_2, \dots, x_n$ avec les probabilités
$p_i = P(X = x_i)$.

\begin{center}
\renewcommand{\arraystretch}{1.5}
\begin{tabular}{|l|c|c|}
  \hline
  & \textbf{Mécanique} & \textbf{Probabilités} \\
  \hline
  Positions
    & $x_i$
    & valeurs prises par $X$ \\
  Poids
    & masses $m_i$ \;($\sum m_i = 1$)
    & probabilités $p_i$ \;($\sum p_i = 1$) \\
  Barycentre
    & $\displaystyle\bar{x} = \sum m_i\, x_i$
    & $\displaystyle E[X] = \sum_i x_i\, P(X = x_i)$ \\
  \hline
\end{tabular}
\end{center}

\needspace{5\baselineskip}
\begin{mdframed}[backgroundcolor=ThemeLight, linecolor=Theme,
  linewidth=1pt, innertopmargin=8pt, innerbottommargin=8pt,
  nobreak=true]
\textbf{Définition.}
L'\textbf{espérance} d'une variable aléatoire discrète $X$ est
\[
  \boxed{E[X] = \sum_{i} x_i \, P(X = x_i).}
\]
C'est le \textbf{centre de gravité} de la loi de $X$.
\end{mdframed}

\medskip

\textbf{Application.}
Soit $X$ le résultat d'un lancer de dé équilibré à~$6$ faces.

\atrou[2.5cm]{%
\[
  E[X]
  = 1 \times \tfrac16 + 2 \times \tfrac16
  + 3 \times \tfrac16 + 4 \times \tfrac16
  + 5 \times \tfrac16 + 6 \times \tfrac16
  = \frac{21}{6} = 3{,}5.
\]
L'espérance n'est pas forcément une valeur possible de $X$.
}

% ====================================================================
\section*{3. Du moment d'inertie à la variance}
% ====================================================================

Revenons à notre poutre.
On sait \emph{où} se trouve le centre de gravité,
mais les charges sont-elles \textbf{regroupées}
autour de ce point ou \textbf{dispersées}
le long de la poutre ?

\medskip

En mécanique, cette dispersion est mesurée par le
\textbf{moment d'inertie} par rapport au centre de gravité :
\[
  I_G = \sum_{i} m_i \,(x_i - \bar{x})^2.
\]

En probabilités, la même idée donne :

\begin{mdframed}[backgroundcolor=ThemeLight, linecolor=Theme,
  linewidth=1pt, innertopmargin=8pt, innerbottommargin=8pt]
\textbf{Définition.}
La \textbf{variance} de $X$ est
\[
  \boxed{\text{Var}(X)
  = E\!\Big[\,(X - E[X])^2\,\Big]
  = \sum_i \bigl(x_i - E[X]\bigr)^2 \, P(X = x_i).}
\]
L'\textbf{écart-type} est $\sigma(X) = \sqrt{\text{Var}(X)}$.
\end{mdframed}

\medskip

\textbf{Application.}
Reprenons notre poutre, avec les poids normalisés
$0{,}2$, $0{,}5$, $0{,}3$ aux positions $1$, $3$, $7$
($\bar{x} = 3{,}8$).

\atrou[3.6cm]{%
\begin{align*}
  \text{Var}
  &= 0{,}2 \times (1-3{,}8)^2
   + 0{,}5 \times (3-3{,}8)^2
   + 0{,}3 \times (7-3{,}8)^2 \\
  &= 0{,}2 \times 7{,}84
   + 0{,}5 \times 0{,}64
   + 0{,}3 \times 10{,}24 \\
  &= 1{,}568 + 0{,}32 + 3{,}072 = 4{,}96.
\end{align*}
}

% ====================================================================
\section*{4. Le théorème de Huygens}
% ====================================================================

En mécanique, le \textbf{théorème de Huygens}
permet de calculer le moment d'inertie
sans repasser par le barycentre :
\[
  I_G = \sum m_i \, x_i^2 \;-\; M\,\bar{x}^{\,2}
  \qquad (M = \text{masse totale}).
\]

Avec $M = 1$ (probabilités), on obtient la formule pratique :

\atrou[4cm]{%
\[
  \boxed{\text{Var}(X) = E[X^2] - \bigl(E[X]\bigr)^2.}
\]
\textbf{Preuve.} On développe le carré dans la définition :
\begin{align*}
  \text{Var}(X)
  &= E\bigl[(X-E[X])^2\bigr]
   = E\bigl[X^2 - 2X\,E[X] + E[X]^2\bigr] \\
  &= E[X^2] - 2\,E[X]\cdot E[X] + E[X]^2
   = E[X^2] - E[X]^2.
\end{align*}
}

\textbf{Vérification sur l'exemple.}

\atrou[2.5cm]{%
$E[X^2] = 0{,}2 \times 1^2 + 0{,}5 \times 3^2
          + 0{,}3 \times 7^2
= 0{,}2 + 4{,}5 + 14{,}7 = 19{,}4.$

$\text{Var}(X) = 19{,}4 - 3{,}8^2 = 19{,}4 - 14{,}44
= 4{,}96.\;\checkmark$
}

% ====================================================================
\section*{5. Propriétés de l'espérance et de la variance}
% ====================================================================

\textbf{Linéarité de l'espérance.}
Pour $a,b \in \mathbb{R}$ :
\[
  E[aX+b] = a\,E[X] + b.
\]

\textbf{Additivité}
(toujours vraie, même sans indépendance) :
\[
  E[X+Y] = E[X] + E[Y].
\]

\medskip

\textbf{Variance et transformation affine :}

\atrou[2cm]{%
\[
  \text{Var}(aX+b) = a^2\,\text{Var}(X).
\]
Translater ne change pas la dispersion ;
dilater multiplie la variance par~$a^2$.
}

\textbf{Variance d'une somme}
(sous hypothèse d'indépendance) :

Si $X$ et $Y$ sont \textbf{indépendantes},
$\;\text{Var}(X+Y) = \text{Var}(X) + \text{Var}(Y).$

\smallskip

\textit{Interprétation mécanique :
si deux systèmes de masses indépendants sont réunis,
les moments d'inertie s'additionnent.}

% ====================================================================
\section*{6. Espérance et variance des lois classiques}
% ====================================================================

Considérons une seule situation de chantier :

\begin{center}
\begin{mdframed}[backgroundcolor=ThemeLight, linecolor=Theme,
  linewidth=1pt, innertopmargin=8pt, innerbottommargin=8pt]
\centering\itshape
Un contrôleur qualité inspecte des pièces sur un chantier.\\
Chaque pièce, indépendamment, a une probabilité $p$ d'être défectueuse.
\end{mdframed}
\end{center}

Selon la \textbf{question posée}, on obtient une loi différente :

\begin{center}
\renewcommand{\arraystretch}{1.4}
\begin{tabular}{|m{6.5cm}|l|}
  \hline
  \textbf{Question} & \textbf{Loi} \\
  \hline
  \og Cette pièce est-elle défectueuse ?\fg{}
    & Bernoulli \\
  \hline
  \og Sur $n$ pièces, combien sont défectueuses ?\fg{}
    & Binomiale \\
  \hline
  \og Quand apparaît le premier défaut ?\fg{}
    & Géométrique \\
  \hline
  \og Combien de défauts sur une très longue série ?\fg{}
    & Poisson \\
  \hline
\end{tabular}
\end{center}

\subsection*{Loi de Bernoulli $\mathcal{B}(p)$ --- \og est-ce défectueux ?\fg{}}

Un seul essai, deux issues.
$X$ vaut $1$ (défaut) avec probabilité $p$
et $0$ (conforme) avec probabilité $1-p$.

\atrou[3cm]{%
\[
  E[X] = 0 \cdot (1-p) + 1 \cdot p = p.
\]
\[
  E[X^2] = 0^2 \cdot (1-p) + 1^2 \cdot p = p
  \;\;\implies\;\;
  \text{Var}(X) = p - p^2 = p(1-p).
\]
}

\subsection*{Loi binomiale $\mathcal{B}(n,p)$ --- \og combien sur $n$ pièces ?\fg{}}

On répète $n$ essais de Bernoulli indépendants
et on \textbf{compte} le nombre de succès.

\medskip

\textbf{Idée clé :}
$X = X_1 + X_2 + \cdots + X_n$
avec $X_i \sim \mathcal{B}(p)$ indépendantes.

\atrou[3cm]{%
Par linéarité de l'espérance :
$\;E[X] = E[X_1] + \cdots + E[X_n] = np.$

\medskip

Par additivité des variances (indépendance) :
$\;\text{Var}(X) = \text{Var}(X_1) + \cdots + \text{Var}(X_n)
= np(1-p).$
}

\subsection*{Loi géométrique $\mathcal{G}(p)$
  --- \og quand le premier défaut ?\fg{}}

On ne compte plus \emph{combien} de défauts,
mais \emph{quand} apparaît le premier.
$T$ est le rang du premier défaut :

\medskip

\textbf{Exemple.}
Les 4 premières pièces sont conformes,
la 5\up{e} est défectueuse : $T = 5$.

\medskip

Pour que $T = k$, il faut $k-1$ conformes puis $1$ défaut :

\atrou[2.5cm]{%
\[
  P(T = k) = \underbrace{(1-p)^{k-1}}_{\text{$k-1$ conformes}}
  \times\; p,
  \qquad k \geq 1.
\]
}

\textbf{Espérance} (admis, démontré en TD) :

\atrou[2cm]{%
\[
  E[T] = \frac{1}{p},
  \qquad
  \text{Var}(T) = \frac{1-p}{p^2}.
\]
\textit{Si $p = 5\,\%$, on attend en moyenne $1/0{,}05 = 20$ pièces
avant le premier défaut.}
}

\textbf{Question.} Les 10 premières pièces sont conformes.
Combien de pièces supplémentaires faut-il attendre
en moyenne avant le premier défaut ?

\atrou[2.5cm]{%
\textbf{Réponse :} encore $1/p = 20$ pièces.
Le passé n'influence pas l'attente future :
c'est la \textbf{propriété sans mémoire}
de la loi géométrique.
(Preuve : $P(T > n+k \mid T > n) = P(T > k)$.)
}

\subsection*{Loi de Poisson $\mathcal{P}(\lambda)$
  --- \og combien sur une grande série ?\fg{}}

Quand $n$ est très grand et $p$ très petit (événements rares),
la loi $\mathcal{B}(n,p)$ est difficile à calculer.
On la remplace par un modèle plus simple :
la loi de Poisson de paramètre $\lambda = np$.

\medskip

$P(X=k)=e^{-\lambda}\dfrac{\lambda^k}{k!}$, \; $k\in\mathbb{N}$.

\[
  E[X] = \lambda,
  \qquad
  \text{Var}(X) = \lambda.
\]

\textit{Propriété remarquable : pour une Poisson,
moyenne et variance sont égales.}

\medskip

\textbf{D'où vient cette formule ?}

\atrou[3cm]{%
On divise une grande zone en $n$ petits morceaux,
chacun ayant une probabilité $\lambda/n$ de contenir un défaut.
Le nombre total suit $\mathcal{B}(n, \lambda/n)$.
Quand $n \to \infty$, on obtient la loi de Poisson.

\textit{(Démonstration complète en TD.)}
}

\subsection*{Tableau récapitulatif}

\begin{center}
\renewcommand{\arraystretch}{1.8}
\begin{tabular}{|l|c|c|c|}
  \hline
  \textbf{Loi} & \textbf{Paramètre(s)}
    & $\boldsymbol{E[X]}$
    & \textbf{Var}$\boldsymbol{(X)}$ \\
  \hline
  Bernoulli $\mathcal{B}(p)$
    & $p \in [0,1]$
    & \val{$p$} & \val{$p(1-p)$} \\
  \hline
  Binomiale $\mathcal{B}(n,p)$
    & $n \in \mathbb{N}^*,\; p \in [0,1]$
    & \val{$np$} & \val{$np(1-p)$} \\
  \hline
  Géométrique $\mathcal{G}(p)$
    & $p \in \,]0,1]$
    & \val{$\dfrac{1}{p}$}
    & \val{$\dfrac{1-p}{p^2}$} \\
  \hline
  Poisson $\mathcal{P}(\lambda)$
    & $\lambda > 0$
    & \val{$\lambda$}
    & \val{$\lambda$} \\
  \hline
\end{tabular}
\end{center}

% ====================================================================
\section*{7. Questions ouvertes pour le TD}
% ====================================================================

\begin{enumerate}
  \item \textbf{Attente du $r$-ème défaut.}
    La loi géométrique donne le rang du premier défaut.
    Et si le chantier s'arrêtait au 3\up{e} défaut,
    combien de pièces faudrait-il inspecter ?
    Quelle loi intervient ?

  \item \textbf{D'où vient la loi de Poisson ?}
    On a affirmé que $\mathcal{B}(n, \lambda/n)$
    tend vers $\mathcal{P}(\lambda)$ quand $n \to \infty$.
    Peut-on le démontrer ?

  \item \textbf{Surbooking aérien.}
    Une compagnie vend $176$ billets pour $174$ places.
    Chaque passager ne se présente pas avec probabilité $3\,\%$.
    La stratégie est-elle rentable ?
\end{enumerate}

% ====================================================================
\section*{8. Bilan}
% ====================================================================

\atrou[6cm]{%
\begin{enumerate}
  \item L'espérance est un \textbf{centre de gravité} :
    $\displaystyle E[X] = \sum_i x_i\,P(X=x_i).$
  \item La variance est un \textbf{moment d'inertie}
    qui mesure la dispersion :
    $\text{Var}(X) = E[X^2] - E[X]^2.$
  \item La \textbf{linéarité} de l'espérance
    et l'\textbf{additivité} des variances (sous indépendance)
    permettent de retrouver $E$ et $\text{Var}$
    des lois composées
    (binomiale = somme de Bernoulli).
  \item Une même situation de chantier produit
    des lois différentes selon la \textbf{question posée} :
    combien ? $\to$ binomiale ;
    quand ? $\to$ géométrique ;
    combien sur une grande échelle ? $\to$ Poisson.
\end{enumerate}
}

\end{document}
